\section{Evaluation}
\label{sec:evaluation}
We evaluate three questions: Can the original verification assertion serve as
an instance oracle? Does copied-state checking retain the information required
by later observations? What evidence and cost are associated with a complete
result? The study combines matched contract configurations, independently
checked proof groups, and development regressions. All numerical results below
come from archived runs whose sources, goal interfaces, and hashes are checked
by the paper build.

\subsection{Design and comparison units}
The paired archive contains 102 tool invocations covering 11 contract
configurations from 10 C body families. Each configuration has a valid and
mutated postcondition. CN performs universal verification, Fulminate checks a
concrete fixture, and \tespec{} checks matched logical data with symbolic
addresses. Repeated boundary policies are collapsed when their outcomes agree.
We use CN to examine the effect of modular summaries and Fulminate to compare
instance verdicts. Bennet and Darcy additionally test the translation of the
least-root property. The original formula and its proved finite rewrite are
reported separately.

A complete instance is the primary outcome. We also report phase results,
residual obligations, proof polarity, and replay cost. Archived UNKNOWN results
retain that status. Historical INVALID labels denote empty entry domains;
ERROR labels denote execution or infrastructure diagnostics. The newer VC
interface in Section~\ref{sec:evidence} does not change archived measurements.

\subsection{Checking original mathematical assertions}
% Generated by scripts/summarize-evidence.py; edit the generator.
\begin{table}[t]
\centering\small
\caption{Independent checking of complete original VC groups.}
\label{tab:proofs}
\begin{tabular}{@{}lrcr@{}}
\toprule
Case & VCs & Evidence & Check (s) \\
\midrule
Cell increment & 4 & Proved & 135.57 \\
Square relation & 3 & Proved & 129.53 \\
Relational loop & 44 & Proved & 137.47 \\
Least-root return & 26 & Proved & 132.06 \\
List traversal & 19 & Proved & 140.12 \\
Least-root observation & 27 & Proved & 124.45 \\
Observation mutant & 27 & Refuted VC & 125.86 \\
\bottomrule
\end{tabular}
\par\smallskip\parbox{\columnwidth}{\footnotesize Time includes certificate compilation, assumption audit, and kernel replay. These are separate archived runs; construction time is excluded. The mutant proves one VC's negation and the other 26 original obligations.}
\end{table}

The least-root return fixture fixes $n=13$ and leaves the cell address symbolic.
It executes four loop iterations and generates 26 obligations. Adding the
quantified intermediate observation in Listing~\ref{lst:root} yields a separate
27-obligation instance. Both complete proof groups pass independent checking,
including entry existence and the original safety and functional goals.
The observation instance receives an overall PASS after proof completion.

The pinned Fulminate, Bennet, and Darcy configurations reject the original
quantified postcondition with an upper-bound inference diagnostic. The
proved finite formulation in Equation~\eqref{eq:finite-root} passes Fulminate
on the same fixture. This experiment demonstrates the ability to retain the
original assertion and locates the alternative's extra task: constructing an
executable reformulation and justifying its equivalence. The generator attempts
for the rewritten formula and the unresolved CN intermediate-annotation
translation remain in the archive; the successful comparative result concerns
the postcondition.

\paragraph{An assertion mutation with unchanged computation.}
The negative instance replaces $\mathsf{LeastRoot}(n,r)$ with
$\mathsf{LeastRoot}(n,r-1)$ before the write. Both variants execute four iterations,
visit five invariant states, write 4, and return 4. The negative proof checks
the negation of the complete assertion VC, including its generated alternatives,
and proves the other 26 obligations. Its recorded result is a checked VC
refutation. The archive has no admitted execution-counterexample certificate
for that proof, so the evidence stops at that scope. This pair tests the
assertion independently of the final output.

Table~\ref{tab:proofs} also includes cell, square-relation, relational-loop,
and list-traversal instances. The seven groups contain six positive groups and
one refuted-VC group. They check original interfaces of 3--44 obligations;
certificate compilation, assumption audit, and kernel replay take
124.45--140.12 seconds per group. These timings exclude proof construction.
Proofs were obtained during development using supplied lemmas and exploratory
proof work; agent success rates are a separate measurement.

\subsection{Common-fragment results and execution preservation}
% Generated by scripts/summarize-evidence.py; edit the generator.
\begin{table}[t]
\centering\small
\caption{Paired subjects: valid / mutated contract outcomes.}
\label{tab:paired}
\begin{tabular}{lccc}
\toprule
Subject & CN & Fulminate & TeSpec \\
\midrule
alias\_shared & V / U & P / F & P / U \\
array\_cells & V / U & P / F & P / U \\
heap\_loop & V / U & P / F & P / U \\
list\_length & V / U & P / F & U / U \\
nested\_call & V / U & P / F & P / U \\
pointer\_cell & V / U & P / F & P / U \\
recursive\_heap & V / U & P / F & P / U \\
recursive\_heap\_weak & U / U & P / F & P / U \\
single\_cell & V / U & P / F & P / U \\
struct\_fields & V / U & P / F & P / U \\
swap & V / U & P / F & P / U \\
\bottomrule
\end{tabular}
\par\smallskip\parbox{\columnwidth}{\footnotesize V: verified; P: pass; F: runtime violation; U: unproved or unknown. CN checks full domains, Fulminate concrete fixtures, and TeSpec bound symbolic domains. The 11 contract configurations cover 10 body families; repeated modes are collapsed.}
\end{table}

On the paired common-fragment configurations, automatic \tespec{} checking
accepts 10 of 11 valid contracts and leaves all 11 mutations UNKNOWN.
Fulminate accepts all 11 valid fixtures and reports a runtime violation for
every mutant. The list configuration requires residual proof work in \tespec{};
a separate 19-obligation traversal certificate is reported in
Table~\ref{tab:proofs}. The paired automatic row retains its original outcome.
Thus direct runtime checking provides more complete automatic negative feedback
on this fragment. The proof-backed route supplies the mathematical-assertion
capability measured above at the cost of residual proof work.

\paragraph{Weak recursive summaries.}
The recursive-cell fixture starts at value 3 with counter 3 and returns the
exact final value 6. Replacing the helper's exact postcondition with an
existential value bounded below by the input leaves actual writes unchanged.
Both \tespec{} and Fulminate retain a passing result. CN's modular verifier
and an unmodified QCP baseline retain the top-level exact-result obligation
because the helper summary is too weak. This probe identifies the effect of
executing the helper body while observing its contract. The exact and weak
rows are two configurations of one body.

\paragraph{Local observations and normalization.}
% Generated by scripts/summarize-evidence.py; edit the generator.
\begin{table}[t]
\centering\small
\caption{Mechanism probes separate complete results from local proof progress.}
\label{tab:mechanisms}
\begin{tabular}{p{.30\columnwidth}p{.59\columnwidth}}
\toprule
Mechanism & Observed change \\
\midrule
Local witness resolution & 12 valid-mode runs: U $\to$ P; 12 incorrect-mode runs: U $\to$ U \\
Checked pure facts & Two configurations: 2 $\to$ 1 residual roots each; 0 new complete passes \\
Full proof completion & Positive least-root: U $\to$ P, all 27 obligations checked \\
VC refutation & Intermediate mutant: checked VC refutation; reachable violation unestablished \\
\bottomrule
\end{tabular}
\end{table}

Resolving local pointer witnesses before choosing the exact local frame changes
12 valid-mode runs from UNKNOWN to PASS; the 12 incorrect-mode runs remain
UNKNOWN. The 48 before/after observations cover 12 cases in two modes.
For two configurations of one imported list instance, checked pure consequences
reduce residual root obligations from two to one. Entry existence remains
unresolved, yielding zero additional complete passes. Cold controller costs
are 129.41 and 130.28 seconds. These probes distinguish successful local
simplification from complete instance acceptance.

\subsection{Engineering coverage and formal evidence}
% Generated by scripts/summarize-evidence.py; edit the generator.
\begin{table}[t]
\centering\small
\caption{Development regression outcomes (1,322 runs).}
\label{tab:regression}
\begin{tabular}{lrl}
\toprule
Verdict & Runs & Interpretation \\
\midrule
PASS & 541 & Recorded checks discharged \\
UNKNOWN & 757 & Unresolved or incomplete \\
ERROR & 14 & Execution/coverage diagnostic \\
INVALID & 10 & Empty entry domain \\
\bottomrule
\end{tabular}
\par\smallskip\parbox{\columnwidth}{\footnotesize All recorded expectations are met. Negative checks include deliberately unsupported cases; expectation agreement is distinct from confirmed violation detection.}
\end{table}

The development suite contains 52 functional programs, each with four inputs,
three variants, and two boundary policies, plus 37 diagnostic scenarios in both
policies: $52\cdot4\cdot3\cdot2+37\cdot2=1{,}322$ checks.
Functional expectations use independent Python oracles cross-checked against
defined native C executions. All recorded expectations are met, with the
outcomes in Table~\ref{tab:regression}. The 779 expected negative checks include
unsupported scenarios; their acceptance criterion permits UNKNOWN. Consequently,
the suite measures avoidance of false acceptance and diagnostic conformance,
while confirmed violations are counted separately. Each body mutation changes
at least one observed outcome relative to its reference.

A broader historical scan covers 267 C files and 1,147 specifications with
5,160 synthetic-input runs: 1,086 PASS, 1,022 INVALID, 2,416 UNKNOWN, and
636 ERROR. Another 287 cases require typed bindings. A later default-strategy
scan has a different denominator: 81 sources, 327 specifications, and 1,404 runs.
It reports 448 PASS (212 execution and 236 specification-only), 726 UNKNOWN,
128 INVALID, and 102 ERROR. These inventories expose integration barriers such
as missing bodies and opaque resources. Their synthetic inputs have no
independent correctness labels. On eight original sort, queue, and reverse
instances, the default workflow produces six UNKNOWN and two ERROR results.
These outcomes locate the remaining execution and proof barriers in imported
libraries.

The current formal replay checks 11 modules and audits 88 principal statements,
including observation erasure, lifting, and coverage. The reference results
support the semantic interfaces in Section~\ref{sec:semantics}. Complete
instance certificates separately check the generated obligations listed in
Table~\ref{tab:proofs}. Extended implementation probes and their unsuccessful
attempts are retained in the supplementary development record.
